\documentclass[12pt,a4paper]{article}

\usepackage[utf8]{inputenc}
\usepackage[T1]{fontenc}
\usepackage[margin=2.5cm]{geometry}
\usepackage{hyperref}
\usepackage{microtype}
\usepackage{parskip}
\usepackage{lmodern}

\hypersetup{colorlinks=true, urlcolor=blue!60!black}

\pagestyle{empty}

\begin{document}

\begin{flushright}
Kundai Farai Sachikonye \\
Germany \\
\href{https://github.com/fullscreen-triangle}{github.com/fullscreen-triangle} \\
\texttt{kundai.sachikonye@bitspark.com} \\[6pt]
15 May 2026
\end{flushright}

\vspace{1em}

Prof.\ Philipp Berens \\
Hertie Institute for AI in Brain Health \\
Department of Data Science \\
University of Tübingen / University Hospital Tübingen

\vspace{1.5em}

\textbf{Re: PhD Position --- ML for Medical Imaging, Neurogenomics and Neuronal Modelling}

\vspace{1em}

Dear Prof.\ Berens,

I am applying for the PhD position in machine learning for medical imaging,
neurogenomics, and neuronal modelling at the Hertie Institute. I completed
my MSc in Computational Biology at the University of Tübingen and am applying
because the three research pillars of this position --- medical imaging, neurogenomics,
and neuronal modelling --- correspond closely to independent research I have been
conducting over the past year, and because Tübingen is where I developed the
computational biology foundation that this work builds on.

On the medical imaging side, I have built
\href{https://github.com/fullscreen-triangle/helicopter}{helicopter}, a multi-modal
microscopy image analysis framework that combines fluorescence, brightfield,
phase-contrast, and spectral modalities through sequential constraint satisfaction.
The architecture avoids the typical dependency on single-modality optical resolution
by treating each physical measurement as an independent exclusion factor, reducing
structural ambiguity from $\sim$10$^{60}$ configurations to unique solutions. Separately,
\href{https://github.com/fullscreen-triangle/lavoisier}{lavoisier}, my mass-spectrometry
metabolomics pipeline, implements a dual-pipeline architecture --- one numerical, one
based on converting spectra into temporal image sequences analysed by a convolutional
neural network (PyTorch) --- that achieves 98.9\% feature extraction accuracy against
NIST reference libraries. Building ML pipelines for biological imaging data where ground
truth is ambiguous and data is structurally heterogeneous is the specific problem these
tools address.

A third imaging context: the BRUT Observatory
(\href{https://remote-photoplethysmography.vercel.app/}{remote-photoplethysmography.vercel.app})
inverts a layered Beer--Lambert optical model of forehead skin to recover
four physiological state parameters from standard RGB camera pixels and the
AC modulation amplitude of the cardiac signal: melanin index (median
$0.5\%$ relative error), total haemoglobin ($10\%$), a vasodilation factor
(exact, from AC amplitude), and a derived skin-temperature estimate via
thermoregulatory linearisation. From the recovered state the observed heart
rate is decomposed as
$\mathrm{HR}_\mathrm{obs} = \mathrm{HR}_\mathrm{int} + \Delta\mathrm{HR}_\mathrm{met}
+ \Delta\mathrm{HR}_{O_2} + \Delta\mathrm{HR}_\mathrm{auto}$,
isolating the genuine sympathovagal drive from thermal ($Q_{10}$-mediated)
and hypoxic (chemoreceptor) contributions. The inverse is algebraic ---
channel separation (blue identifies melanin, red identifies Hb, green
identifies oxygenation trend, AC amplitude identifies vasodilation) ---
and runs in sub-microsecond time. The full pipeline runs in real time
in a WebGPU-enabled browser at $1\,\mathrm{Hz}$ state estimates from a
$720\mathrm{p}$ webcam. The paper is candid: oxygenation from three-band
RGB has a structural error ceiling near $40\%$ and is best treated as a
relative-trend signal rather than a clinical reading. The autonomic residual
$\Delta\mathrm{HR}_\mathrm{auto}$ --- the neural-drive component that
does not arise from temperature or hypoxia --- connects directly to the
neuronal modelling pillar of this position.

On the neurogenomics side, \href{https://github.com/fullscreen-triangle/gospel}{gospel}
is a whole-genome variant analysis framework that processes million-scale VCF files
alongside RNA-seq expression matrices, validated against 1000 Genomes Phase 3,
gnomAD v3.1.2, and UK Biobank. The system integrates Bayesian network inference,
Gaussian mixture models for noise characterisation, and variational Bayes for approximate
inference. For retinal neurogenomics specifically, the genomic variant and expression
analysis capabilities map directly onto the kinds of datasets generated in
ophthalmological research: identifying genetic determinants of retinal disease requires
exactly this kind of joint variant-and-expression pipeline operating at scale.

On the neuronal modelling side, I have developed a formal mathematical framework ---
published as \textit{Syndrome: Categorical Resolution of Disease Trajectories} ---
for representing cellular state through eight oscillator classes, one of which is the
calcium signalling oscillator ($D_{Ca}$). Retinal photoreceptors are paradigmatically
calcium-dependent: the phototransduction cascade is fundamentally a calcium-mediated
signal. The coherence index formalism I have developed ($\eta_{Ca} \in [0,1]$ from
regularity of calcium oscillations) provides a quantitative, computationally tractable
representation of photoreceptor health state that could be derived from
electrophysiological or imaging data. A companion paper ---
\textit{On Disease Epistemology in Blind Receiver} --- proves that loop-holonomy
self-consistency of the biochemical circuit (AUC$=$1.00 vs 0.896 for template comparison
across 25 validation experiments) is the unique viable diagnostic criterion when template
states are unavailable, which is the regime of early retinal disease.

Separately, I have developed a first-principles model of nervous system dynamics as
a closed electrical circuit. The foundational axiom --- $\oint \mathbf{J} \cdot d\mathbf{A} = 0$,
no external charge ground --- derives that centre-of-pressure motion decomposes into
two independent components: \textit{rambling} (supraspinal drift, frequency below a
critical $f_c$) and \textit{trembling} (peripheral oscillation, above $f_c$), corresponding
to cortical and spinal contributions respectively. Validated across 86 consecutive
nights of wearable data at five recording sites and across Parkinson's disease,
cerebellar ataxia, aging, and dual-task conditions, the framework extracts
$Q_\text{motor} = 290.9$\,mC/s and $Q_\text{thought} = 133.5$\,mC/s as empirical
constants from longitudinal data, with a dream-to-thought charge ratio of 0.975
matching the theoretical prediction $\sqrt{0.95}$ from first principles.
The deafferentation paradox --- that sensory deafferentation does not suppress neural
oscillation but maintains it --- is reinterpreted as circuit interruption: the
oscillator continues because it has no external ground to discharge into.
This connects directly to retinal disease: photoreceptor degeneration is
visual deafferentation, and the maintained oscillatory activity in ganglion cells
following photoreceptor loss is precisely what this circuit model predicts.
The live tools are available at
\href{https://rembling-trambling.vercel.app/tools/rambling}{rembling-trambling.vercel.app}.

I have also considered the specific connection to ophthalmological autoimmunity.
My paper on \textit{VDJ Recombination Categorical Filter Cascades} predicts that proteins
with high categorical richness $R$ --- quantifying isoform diversity, conformational entropy,
and downstream connectivity --- are preferentially targeted by the immune system
($r=0.73$, $p<10^{-12}$, $n=2{,}847$ IEDB peptides). Retinal antigens implicated in
autoimmune retinopathy (recoverin, $\alpha$-enolase, CRABP-1) are precisely the kind
of proteins with high conformational flexibility and isoform diversity that would be
predicted to cross the categorical threshold. This is a specific, testable connection to
the group's interest in diseased retinal biology.
The spectral framework underlying the MHC binding prediction is now deployed at
\href{https://syndrome-seven.vercel.app}{syndrome-seven.vercel.app}: the MHC binding
predictor converts a peptide sequence into a 36-dimensional spectral embedding via
3-channel DFT and returns cosine similarity $\rho$ against the MHC groove aperture
in $O(cL\log L)$ without a database query. The neoantigen tool applies a simultaneous
dual threshold --- self-distinction $\Delta\rho_\text{self} > 0.15$ and MHC presentation
$\rho_\text{MHC} > 0.72$ --- to any wild-type / mutant peptide pair, returning whether
the mutation is immunologically visible and presentable. Applied to recoverin peptides,
this provides a direct in-browser test of which specific residues, when mutated or
post-translationally modified, cross the autoimmune threshold.

My programming work is primarily Python and Rust, with R from my bioinformatics training.
I am not an ophthalmologist and would not claim to be: the retinal domain knowledge is
what a doctoral programme provides. I am applying because the ML methodology --- building
pipelines that extract signal from structurally complex, high-dimensional biological
imaging and genomics data, and modelling disease state from multi-modal biological
measurements --- is precisely what I have been building for a year, and because Tübingen
is where I am equipped to do this work well.

I am available immediately, fully authorised to work in Germany (Aufenthaltserlaubnis),
and would welcome the opportunity to discuss this application.

\vspace{1.5em}
Yours sincerely,

\vspace{2.5em}
\textbf{Kundai Farai Sachikonye}

\end{document}
